\documentclass[12pt]{article}

\usepackage[a4paper, margin=2.5cm]{geometry}
\usepackage{amsmath}
\usepackage{amssymb}
\usepackage{amsthm}
\usepackage[colorlinks=true, linkcolor=blue, citecolor=blue, urlcolor=blue]{hyperref}
\usepackage{booktabs}
\usepackage{natbib}

\newtheorem{theorem}{Theorem}
\newtheorem{proposition}[theorem]{Proposition}
\newtheorem{definition}[theorem]{Definition}
\newtheorem{remark}[theorem]{Remark}
\newtheorem{example}[theorem]{Example}

\newcommand{\Mcal}{\mathcal{M}}
\newcommand{\Ccal}{\mathcal{C}}
\newcommand{\Fcal}{\mathcal{F}}
\newcommand{\Pcal}{\mathcal{P}}
\newcommand{\Sadm}{\mathcal{S}_{\mathrm{adm}}}
\newcommand{\Scl}{\mathcal{S}_{\mathrm{cl}}}
\newcommand{\phig}{\varphi}

\title{Confinement as a Multi-Node Closure Constraint\\in $\phig$-Structured Geometry}
\author{%
  Lee Smart\\[4pt]
  \textit{Institute of Vibrational Field Dynamics}\\[2pt]
  \texttt{contact@vibrationalfielddynamics.org}\\[2pt]
  \texttt{@vfd\_org}%
}
\date{April 2026}

\begin{document}

\maketitle

\noindent\textbf{Status:} Preprint (not peer-reviewed)

\begin{abstract}
We extend the closure-constraint framework of Paper~VII~\citep{SmartVII} to composite multi-node systems on the $\phig$-structured geometry. The central result is a reformulation of confinement as a property of the closure constraint structure: states with disconnected shell support are excluded from the closed state set $\Scl$ by the connectivity constraint, and therefore cannot be realised as closure-stable configurations.

Connected composite states --- such as the proton on shells $\{2,3,4\}$ --- satisfy all closure constraints and correspond to fixed points of the closure operator. Disconnected states --- such as a quark-like configuration on $\{2,4\}$ --- violate the connectivity constraint and carry a nonzero gap penalty, placing them outside $\Scl$. Confinement is therefore reinterpreted as the statement that only connected composites are closure-stable; isolated quark-like states are not excluded by a confining force, but by the constraint structure itself.

The paper introduces a composite closure functional that extends the variational formulation of Paper~VII with an explicit gap-penalty term. Emergent composite classes (single-node, two-node, three-node) are identified, with disconnected configurations excluded. The framework is structurally compatible with the observed absence of free quarks and with the composite structure of hadrons, but does not derive the $\mathrm{SU}(3)$ gauge group, gluon dynamics, or asymptotic freedom.
\end{abstract}

%==================================================================
\section{Introduction}\label{sec:intro}
%==================================================================

Paper~VII~\citep{SmartVII} formalised the closure operator, the closed state set, and the constraint classes governing single-node and multi-node configurations on the $\phig$-structured geometry. That paper demonstrated that closure constraints restrict the admissible solution space in a non-trivial way: generic states do not satisfy closure, and the constraint structure follows from the spectral and combinatorial properties of the underlying geometry.

The present paper extends that framework to composite systems. The central physical question is: \textit{why do isolated quark-like states not appear as free particles?}

In the Standard Model, confinement is a dynamical property of the $\mathrm{SU}(3)$ gauge theory: the strong coupling grows at large distances, preventing colour-charged states from being isolated. This paper offers a complementary structural interpretation: confinement as a consequence of multi-node closure constraints. States with disconnected shell support are excluded from $\Scl$ not by a force, but by the constraint structure itself.

\subsection*{Scope and epistemic status}

This paper reformulates confinement as a property of the closure constraint structure rather than as a force mediated by gauge fields. It does not derive the $\mathrm{SU}(3)$ gauge group, compute scattering amplitudes, or reproduce lattice QCD results. The framework is presented as a structural reinterpretation compatible with observed hadronic structure, not as a replacement for quantum chromodynamics.

%==================================================================
\section{Composite State Space}\label{sec:composite}
%==================================================================

\subsection{Single-node states}

In the framework of Paper~VII, a single-node state $\psi \in \Sadm$ is characterised by its shell support $S \subseteq \{1, \ldots, 5\}$, winding number $w$, and torsional configuration. Single-node states describe elementary configurations: the electron on $\{1\}$, leptons on single shells, and the fundamental building blocks of composite structures.

\subsection{Composite states}

A composite state is a multi-component configuration:
\begin{equation}\label{eq:composite}
    \Psi = (\psi_1, \psi_2, \ldots, \psi_n)
\end{equation}
where each component $\psi_i$ occupies a subset of the shell structure. The composite shell support is the union:
\begin{equation}
    S_\Psi = \bigcup_{i=1}^{n} S_i
\end{equation}

The composite state inherits the shell-support graph $G(S_\Psi)$: one vertex per occupied shell, one edge between each pair of consecutive occupied shells. The graph-theoretic properties of $G(S_\Psi)$ --- connectivity, degree sequence, spectral structure --- determine the closure properties of the composite.

\begin{definition}[Composite admissible space]\label{def:compadm}
The composite admissible space $\Sadm^{(n)}$ consists of all normalised $n$-component configurations satisfying the boundary conditions of the physical context. For the 600-cell with $N = 5$ shells, the composite space is finite-dimensional.
\end{definition}

%==================================================================
\section{Multi-Node Closure Constraints}\label{sec:constraints}
%==================================================================

We extend the constraint classes of Paper~VII to composite systems. The key new ingredient is the \textit{connectivity constraint}, which enforces connected shell support.

\subsection{Connectivity constraint}

\begin{definition}[Connectivity]\label{def:connectivity}
A composite state $\Psi$ with shell support $S_\Psi$ satisfies the connectivity constraint if the shell-support graph $G(S_\Psi)$ is connected: there exists a path of consecutive shells linking every pair of occupied shells.
\end{definition}

States with disconnected support --- where $G(S_\Psi)$ has two or more connected components --- violate this constraint.

\subsection{Gap penalty}

\begin{definition}[Gap penalty]\label{def:gap}
The gap penalty $\Xi(\Psi) \geq 0$ measures the degree of disconnection in the composite support:
\begin{equation}\label{eq:gap}
    \Xi(\Psi) = \sum_{\text{gaps}} \left(k_{i+1} - k_i - 1\right)
\end{equation}
where the sum runs over all pairs of consecutive occupied shells $(k_i, k_{i+1})$ with $k_{i+1} - k_i > 1$ (i.e.\ over all gaps in the support). For connected supports, $\Xi = 0$. For disconnected supports, $\Xi > 0$.
\end{definition}

\begin{example}
The proton support $S = \{2,3,4\}$ has no gaps: $\Xi = 0$. A quark-like configuration on $S = \{2,4\}$ has one gap (shell 3 missing): $\Xi = 1$. A configuration on $S = \{1,4\}$ has a two-shell gap: $\Xi = 2$.
\end{example}

\subsection{Extended constraint classes}

The composite closure constraints extend those of Paper~VII:

\begin{itemize}
    \item \textbf{Class I (local):} adjacency consistency, spectral compatibility, phase matching --- as in Paper~VII, applied to each occupied shell.
    \item \textbf{Class II (multi-node):} connectivity ($\Xi = 0$), composite stability ($C(S)$ in admissible range), winding consistency.
    \item \textbf{Class III (projection):} observable equivalence class preserved under admissible torsional deformations --- as in Paper~VII.
\end{itemize}

The connectivity constraint is the new element. It elevates the informal observation from Paper~V~\citep{SmartV} (that the up quark on $\{2,4\}$ is disconnected and must bind) to a formal constraint within the closure framework.

%==================================================================
\section{Composite Closure Functional}\label{sec:functional}
%==================================================================

The composite closure functional extends the variational formulation of Paper~VII with the gap-penalty term:
\begin{equation}\label{eq:compfunctional}
    \Fcal_{\mathrm{comp}}(\Psi) = w_1\, d(\Psi) + w_2\, \mathrm{Var}(\Psi) + w_3\, \tau(\Psi) + w_4\, \Xi(\Psi)
\end{equation}

where:
\begin{itemize}
    \item $d(\Psi)$: local compatibility defect (Class~I violations).
    \item $\mathrm{Var}(\Psi)$: structural instability.
    \item $\tau(\Psi)$: torsional mismatch (Class~III violations).
    \item $\Xi(\Psi)$: gap penalty (Class~II connectivity violation).
\end{itemize}

The weights $w_1, w_2, w_3, w_4 > 0$ are structural parameters, not uniquely determined in this paper.

\begin{remark}
The composite closure functional reduces to the Paper~VII functional when $\Xi = 0$ (i.e.\ for connected configurations). It extends it by adding a penalty for disconnected support, which is the mechanism through which confinement enters the framework.
\end{remark}

A composite state $\Psi$ is in the closed set $\Scl$ if and only if $\Fcal_{\mathrm{comp}}(\Psi) = 0$: all constraint classes are satisfied, including connectivity.

%==================================================================
\section{Confinement as Closure Constraint}\label{sec:confinement}
%==================================================================

\begin{remark}[Confinement as constraint exclusion]\label{rem:confinement}
Within the present framework, confinement is reinterpreted as the statement that states with disconnected shell support are excluded from $\Scl$. Specifically: if $G(S_\Psi)$ is disconnected, then $\Xi(\Psi) > 0$, so $\Fcal_{\mathrm{comp}}(\Psi) > 0$, and therefore $\Psi \notin \Scl$. Such states are not closure-stable and are not realised as physical particles.

Confinement is not imposed by an external force law. It is a consequence of the connectivity constraint within the closure framework: the constraint structure itself excludes isolated quark-like configurations.
\end{remark}

This reinterpretation is structural, not dynamical. It does not replace the dynamical mechanism of confinement in QCD (colour flux tubes, running coupling, area law); it provides a complementary geometric perspective in which confinement appears as a constraint property rather than a force property.

%==================================================================
\section{Worked Examples}\label{sec:examples}
%==================================================================

\subsection{Quark-like state: $S = \{2,4\}$}

\begin{example}[Disconnected quark-like state]\label{ex:quark}
A state with shell support $S = \{2,4\}$ has:
\begin{itemize}
    \item $G(S) = $ two isolated vertices (shells 2 and 4, with shell 3 missing).
    \item Gap penalty: $\Xi = 4 - 2 - 1 = 1 > 0$.
    \item Connectivity: violated.
    \item Closure status: $\Psi \notin \Scl$.
\end{itemize}
This state cannot sustain a standing wave across the gap at shell 3. It must bind with other components whose support fills the gap to form a connected composite.
\end{example}

\subsection{Proton: $S = \{2,3,4\}$}

\begin{example}[Connected baryon]\label{ex:proton}
The proton configuration on $S = \{2,3,4\}$ has:
\begin{itemize}
    \item $G(S) = P_3$ (path graph on 3 vertices): connected.
    \item Gap penalty: $\Xi = 0$.
    \item Closure invariant: $C(\{2,3,4\}) = 24 - 9 - 1 = 14$; $\Delta C = 15$.
    \item Closure status: $\Psi \in \Scl$ (all constraints satisfied).
    \item Mass: $m_p/m_e = \phig^{1265/81} \approx 1835.8$ (Paper~IV~\citep{SmartIV}).
\end{itemize}
The proton fills the connectivity gap of the $\{2,4\}$ quark-like state. It is a closure-stable attractor: a fixed point of the composite closure operator.
\end{example}

\subsection{Neutron}

\begin{example}[Neutron as closure perturbation]\label{ex:neutron}
The neutron shares the proton's shell support $\{2,3,4\}$ and is closure-stable. It differs from the proton in the symmetry sector (charged vs.\ neutral baryon), corresponding to a small perturbation in the torsional configuration. The neutron-proton mass splitting ($\Delta m \approx 1.3$ MeV) is compatible with a small closure-residual difference within the same attractor basin, but is not derived in the present framework.
\end{example}

%==================================================================
\section{Emergent Composite Classes}\label{sec:classes}
%==================================================================

The connectivity constraint, combined with the shell structure of the 600-cell ($N = 5$ shells), induces a classification of composite configurations:

\begin{table}[ht]
\centering
\caption{Emergent composite classes from closure constraints.}
\label{tab:classes}
\begin{tabular}{@{}llll@{}}
\toprule
Class & Shell support type & Connectivity & Physical interpretation \\
\midrule
Single-node & $|S| = 1$ & Trivially connected & Leptons, reference states \\
Two-node (connected) & $|S| = 2$, consecutive & Connected & Meson-like composites \\
Three-node (connected) & $|S| = 3$, consecutive & Connected & Baryon-like composites \\
Disconnected & Any $S$ with gaps & Violated & Excluded from $\Scl$ \\
\bottomrule
\end{tabular}
\end{table}

\begin{remark}
The three-node connected class includes the proton and neutron on $\{2,3,4\}$. The assignment rules (Paper~V~\citep{SmartV}) further constrain which specific supports within each class are admissible. The classification here is purely structural: it follows from connectivity and the $N = 5$ shell geometry.
\end{remark}

%==================================================================
\section{Colour-Like Degrees of Freedom}\label{sec:colour}
%==================================================================

The Standard Model assigns quarks a three-valued colour charge under $\mathrm{SU}(3)$, with confinement requiring colour-neutral composites. Within the present framework, we do not derive $\mathrm{SU}(3)$ or the colour quantum number. However, we note a structural parallel.

The baryon-like composites on $\{2,3,4\}$ involve three distinct shells, each carrying its own local Hilbert space and torsional structure. The requirement that the composite be closure-stable imposes constraints on how these three components combine --- a threefold composite admissibility condition that is structurally reminiscent of the colour-singlet requirement.

\begin{remark}
This observation is a structural analogy, not a derivation of $\mathrm{SU}(3)$. The threefold structure arises here from the geometry (three consecutive interior shells in a five-shell manifold), not from a gauge principle. Whether this geometric threefold structure can be formally related to the $\mathrm{SU}(3)$ colour group is an open question.
\end{remark}

%==================================================================
\section{Stability and Attractors}\label{sec:attractors}
%==================================================================

The composite closure operator $\Ccal_{\mathrm{comp}}$ acts on $\Sadm^{(n)}$ by projecting onto the composite closed set:
\begin{equation}
    \Ccal_{\mathrm{comp}}(\Psi) \in \arg\inf_{\Phi \in \Scl^{(n)}} \|\Psi - \Phi\|
\end{equation}

Fixed points of $\Ccal_{\mathrm{comp}}$ are exactly the states in $\Scl^{(n)}$: composite configurations satisfying all closure constraints, including connectivity.

The attractor basins of the composite closure operator provide the framework-level notion of which initial configurations are associated with which composite state type, whenever the iterative relaxation converges:
\begin{equation}
    B(\Psi^*) = \left\{\Psi_0 \in \Sadm^{(n)} \;\middle|\; \lim_{k \to \infty} \Ccal_{\mathrm{comp}}^k(\Psi_0) = \Psi^*\right\}
\end{equation}

Distinct attractor basins correspond to distinct composite state classes. The proton and neutron occupy the same connectivity class ($\{2,3,4\}$) but different torsional/symmetry sectors, and may be interpreted as distinct attractors within the same structural neighbourhood.

%==================================================================
\section{Relation to QCD}\label{sec:qcd}
%==================================================================

The present framework is structurally compatible with the following features of QCD:

\begin{itemize}
    \item \textbf{Confinement:} isolated quark-like states are excluded from $\Scl$ by the connectivity constraint.
    \item \textbf{Absence of free quarks:} disconnected configurations carry $\Xi > 0$ and are not closure-stable.
    \item \textbf{Composite structure:} baryons correspond to three-node connected composites; mesons to two-node connected composites.
    \item \textbf{Threefold multiplicity:} the three interior shells provide a structural parallel to the three colour degrees of freedom.
\end{itemize}

The framework does not reproduce:

\begin{itemize}
    \item The $\mathrm{SU}(3)$ gauge group or gluon dynamics.
    \item Asymptotic freedom or the running of the strong coupling.
    \item Scattering amplitudes or cross-sections.
    \item Lattice QCD results or the string tension.
\end{itemize}

\textsc{Status:} The confinement reinterpretation is a structural observation within the closure framework. It demonstrates compatibility with observed hadronic structure but does not constitute a derivation of QCD from the $\phig$-structured geometry.

%==================================================================
\section{Limitations}\label{sec:limits}
%==================================================================

\begin{itemize}
    \item \textbf{$\mathrm{SU}(3)$ is not derived.} The colour gauge group is not predicted by the framework. The threefold structure is a geometric observation, not a gauge-theoretic derivation.
    \item \textbf{No dynamics.} The framework addresses static constraint structure, not dynamical evolution, scattering, or decay.
    \item \textbf{No running coupling.} Asymptotic freedom and infrared slavery are dynamical properties not captured by the static closure framework.
    \item \textbf{No lattice comparison.} The framework does not produce predictions comparable to lattice QCD computations.
    \item \textbf{Constraint weights are not uniquely fixed.} The composite closure functional includes structural parameters $w_1, \ldots, w_4$ that are not derived.
    \item \textbf{The gap penalty is a modelling choice.} The specific form of $\Xi(\Psi)$ (Eq.~\ref{eq:gap}) is one of several possible definitions. The qualitative result (disconnected states excluded) is robust to the choice of penalty, but the quantitative details depend on it.
\end{itemize}

%==================================================================
\section{Conclusion}\label{sec:conclusion}
%==================================================================

Confinement is reinterpreted within the closure framework as a consequence of multi-node connectivity constraints. States with disconnected shell support are excluded from the closed state set $\Scl$ by the gap penalty; only connected composite configurations are closure-stable. The proton on $\{2,3,4\}$ is a closure-stable attractor; the quark-like configuration on $\{2,4\}$ is not, because it violates connectivity.

The composite closure functional extends the variational formulation of Paper~VII with an explicit gap-penalty term, providing a quantitative measure of disconnection. Emergent composite classes (single-node, two-node, three-node) follow from the connectivity constraint and the five-shell geometry. The threefold structure of baryon-like composites provides a structural parallel to the colour degrees of freedom, though $\mathrm{SU}(3)$ is not derived.

The framework is structurally compatible with the observed absence of free quarks and the composite structure of hadrons. It provides a geometric perspective in which confinement appears as a constraint property of the $\phig$-structured geometry rather than as a force mediated by gauge fields. Extension to dynamical properties, gauge-group derivation, and comparison with lattice QCD are left for future work.

%==================================================================
\bibliographystyle{plainnat}
\bibliography{references}

\end{document}
